\documentclass[12pt,a4paper]{report}

\usepackage[utf8]{inputenc}
\usepackage[T1]{fontenc}
\usepackage[french]{babel}
\usepackage{geometry}
\usepackage{setspace}
\usepackage{xcolor}
\usepackage{hyperref}
\usepackage{graphicx}
\usepackage{booktabs}
\usepackage{longtable}
\usepackage{array}
\usepackage{float}
\usepackage{caption}
\usepackage{enumitem}
\usepackage{listings}
\usepackage{csquotes}
\usepackage[backend=biber,style=ieee]{biblatex}

\addbibresource{bibliographie.bib}

\geometry{left=2.8cm,right=2.4cm,top=2.5cm,bottom=2.5cm}
\onehalfspacing
\hypersetup{colorlinks=true,linkcolor=blue!50!black,urlcolor=blue!50!black,citecolor=blue!50!black}

\definecolor{sfgreen}{RGB}{38,120,75}
\definecolor{sfbrown}{RGB}{106,76,42}
\definecolor{sfgray}{RGB}{245,247,250}
\definecolor{codegray}{RGB}{248,248,248}

\lstdefinelanguage{mermaid}{
  morekeywords={graph,flowchart,sequenceDiagram,classDiagram,erDiagram,stateDiagram-v2,participant,actor,as,subgraph,end,loop,alt,opt,rect,note,over,activate,deactivate},
  sensitive=false,
  morecomment=[l]{\%\%}
}

\lstset{
  basicstyle=\ttfamily\small,
  breaklines=true,
  frame=single,
  backgroundcolor=\color{codegray},
  columns=fullflexible,
  keepspaces=true
}

\graphicspath{{figures/}}

\newcommand{\screenshotfigure}[3]{
\begin{figure}[H]
\centering
\IfFileExists{figures/#1}{
  \includegraphics[width=0.92\textwidth]{#1}
}{
  \fbox{
    \begin{minipage}[c][6cm][c]{0.88\textwidth}
    \centering
    \textbf{Capture d'écran à insérer}\\[0.4cm]
    Fichier attendu : \texttt{figures/#1}\\[0.2cm]
    #2
    \end{minipage}
  }
}
\caption{#2}
\label{#3}
\end{figure}
}

\begin{document}

\begin{titlepage}
\centering
\vspace*{1cm}
{\Large République Tunisienne}\\[0.2cm]
{\large Ministère de l'Enseignement Supérieur et de la Recherche Scientifique}\\[1.5cm]
{\Huge\bfseries Rapport de Projet de Fin d'Etudes}\\[0.8cm]
{\LARGE\bfseries Smart Firma}\\[0.3cm]
{\Large Plateforme intelligente d'aide à la décision agricole}\\[1cm]
\rule{0.85\textwidth}{0.5pt}\\[0.5cm]
{\large Technologies principales : Django REST Framework, React, TypeScript, SQLite, TensorFlow/Keras, Scikit-learn, Open-Meteo, NewsAPI, OpenAI API}\\[0.5cm]
\rule{0.85\textwidth}{0.5pt}\\[1.5cm]
\begin{tabular}{ll}
Réalisé par : & \textit{Nom de l'étudiant} \\
Encadrant académique : & \textit{Nom de l'encadrant} \\
Organisme d'accueil : & \textit{Smart Firma / organisme à compléter} \\
Année universitaire : & 2025--2026 \\
\end{tabular}
\vfill
{\large Mai 2026}
\end{titlepage}

\pagenumbering{roman}

\chapter*{Remerciements}
\addcontentsline{toc}{chapter}{Remerciements}
Je tiens à exprimer ma profonde reconnaissance à mon encadrant académique pour ses orientations, ses remarques et son suivi tout au long de la réalisation de ce projet. Je remercie également l'équipe pédagogique pour la qualité de la formation reçue, ainsi que toutes les personnes ayant contribué directement ou indirectement à l'aboutissement de ce travail. Enfin, mes remerciements s'adressent aux utilisateurs et aux acteurs du domaine agricole dont les besoins ont inspiré la conception fonctionnelle de la plateforme Smart Firma.

\chapter*{Dédicaces}
\addcontentsline{toc}{chapter}{Dédicaces}
Je dédie ce travail à ma famille, pour son soutien moral et ses encouragements continus, ainsi qu'à toutes les personnes qui m'ont accompagné durant mon parcours académique. Ce projet représente l'aboutissement d'un effort de recherche, de développement et d'apprentissage autour des technologies web et de l'intelligence artificielle appliquées à l'agriculture.

\chapter*{Résumé}
\addcontentsline{toc}{chapter}{Résumé}
Smart Firma est une plateforme web intelligente destinée à accompagner les agriculteurs tunisiens dans leurs décisions quotidiennes. L'application regroupe plusieurs services : authentification, tableaux de bord, météo agricole, actualités, messagerie, chatbot spécialisé, recommandation de cultures et diagnostic de maladies végétales par image. Le frontend est développé avec React et TypeScript, tandis que le backend repose sur Django REST Framework. La solution intègre deux briques d'intelligence artificielle : un modèle Random Forest pour recommander les cultures selon les variables agro-climatiques et un réseau de neurones convolutif Keras pour classifier les maladies des plantes. Le rapport présente l'analyse du code source, l'architecture logicielle, les APIs REST, la base de données, la sécurité, la démarche SCRUM et la méthodologie CRISP-DM appliquée aux modules Machine Learning.

\textbf{Mots-clés :} agriculture intelligente, Django, React, API REST, SCRUM, CRISP-DM, Random Forest, CNN, TensorFlow, Scikit-learn.

\chapter*{Abstract}
\addcontentsline{toc}{chapter}{Abstract}
Smart Firma is an intelligent web platform designed to support Tunisian farmers in daily decision-making. The application provides authentication, dashboards, agricultural weather services, news, messaging, a specialized chatbot, crop recommendation and image-based plant disease diagnosis. The frontend is implemented with React and TypeScript, while the backend relies on Django REST Framework. The system integrates two artificial intelligence components: a Random Forest model for crop recommendation based on agro-climatic features and a Keras convolutional neural network for plant disease classification. This report analyses the source code, software architecture, REST APIs, database design, security mechanisms, SCRUM workflow and the CRISP-DM methodology applied to the Machine Learning modules.

\textbf{Keywords:} smart agriculture, Django, React, REST API, SCRUM, CRISP-DM, Random Forest, CNN, TensorFlow, Scikit-learn.

\chapter*{Liste des abréviations}
\addcontentsline{toc}{chapter}{Liste des abréviations}
\begin{longtable}{p{3cm}p{11cm}}
API & Application Programming Interface \\
CNN & Convolutional Neural Network \\
CRUD & Create, Read, Update, Delete \\
CRISP-DM & Cross Industry Standard Process for Data Mining \\
DRF & Django REST Framework \\
IA & Intelligence Artificielle \\
JWT & JSON Web Token \\
ML & Machine Learning \\
MCD & Modèle Conceptuel de Données \\
MLD & Modèle Logique de Données \\
NPK & Azote, Phosphore, Potassium \\
REST & Representational State Transfer \\
UML & Unified Modeling Language \\
\end{longtable}

\tableofcontents
\listoffigures
\listoftables
\clearpage
\pagenumbering{arabic}

\chapter{Cadre général du projet}

\section{Présentation de l'organisme d'accueil}
Le projet Smart Firma s'inscrit dans un contexte de transformation numérique du secteur agricole. L'organisme d'accueil n'étant pas explicitement décrit dans le code source, le rapport considère Smart Firma comme une structure projet ou une plateforme applicative dédiée à l'agriculture intelligente. Les informations administratives exactes de l'entreprise ou de l'établissement d'accueil doivent être complétées par l'étudiant.

\subsection{Présentation générale}
Smart Firma vise à fournir un environnement numérique unifié permettant aux agriculteurs de consulter des informations utiles, d'obtenir des recommandations et d'exploiter l'intelligence artificielle pour améliorer la surveillance des cultures. La plateforme cible particulièrement le contexte tunisien, avec des cultures telles que l'olivier, le blé dur, les dattes, les agrumes, l'orge, les fèves, le pois chiche, l'amandier, le figuier de barbarie et le pêcher.

\subsection{Activités et services}
Les services détectés dans le code source sont :
\begin{itemize}[noitemsep]
  \item un service d'authentification et de gestion du profil utilisateur ;
  \item un tableau de bord utilisateur et un tableau de bord administrateur ;
  \item un service météo basé sur Open-Meteo et, en secours, OpenWeatherMap ;
  \item un service d'actualités utilisant NewsAPI ;
  \item un chatbot agricole basé sur l'API OpenAI ;
  \item un module de messagerie interne ;
  \item un module de recommandation de cultures tunisiennes ;
  \item un module de diagnostic de maladies végétales par image.
\end{itemize}

\subsection{Organigramme fonctionnel}
Le code ne contient pas d'organigramme officiel de l'organisme. L'organisation fonctionnelle du projet peut néanmoins être représentée ainsi :

\begin{figure}[H]
\centering
\caption{Organigramme fonctionnel proposé pour Smart Firma}
\begin{lstlisting}[language=mermaid]
flowchart TD
    Direction[Direction / Encadrement]
    PO[Product Owner]
    Dev[Equipe de developpement]
    IA[Equipe IA / Data]
    Users[Utilisateurs agricoles]

    Direction --> PO
    PO --> Dev
    PO --> IA
    Dev --> Frontend[Frontend React]
    Dev --> Backend[Backend Django REST]
    IA --> ML[Modeles ML et DL]
    Users --> PO
\end{lstlisting}
\end{figure}

\section{Contexte du projet}
\subsection{Problématique}
Les agriculteurs doivent prendre des décisions rapides concernant les cultures, l'irrigation, les risques météorologiques et les maladies végétales. Les informations nécessaires sont souvent dispersées entre différentes sources : bulletins météo, conseils agronomiques, observations terrain, analyses de sol et connaissances empiriques. Cette dispersion complique la prise de décision et peut entraîner des pertes de rendement ou un usage non optimal des ressources.

\subsection{Objectifs}
Les objectifs du projet sont :
\begin{itemize}[noitemsep]
  \item centraliser les services numériques agricoles dans une plateforme web ;
  \item sécuriser l'accès via inscription, connexion et gestion de session par token ;
  \item fournir des prévisions météorologiques exploitables pour l'agriculture ;
  \item recommander des cultures selon les caractéristiques du sol, du climat et de la localisation ;
  \item diagnostiquer les maladies végétales à partir d'images ;
  \item conserver les historiques d'analyse et de conversation ;
  \item proposer un chatbot spécialisé dans l'agriculture tunisienne ;
  \item documenter la solution selon une démarche académique et professionnelle.
\end{itemize}

\subsection{Motivation}
La motivation principale est d'appliquer les technologies web modernes et l'intelligence artificielle à un domaine à fort impact : l'agriculture. Dans le contexte tunisien, les ressources hydriques, la diversité des zones agro-écologiques et les risques phytosanitaires justifient l'utilisation d'outils numériques d'aide à la décision.

\section{Analyse de l'existant}
\subsection{Solutions concurrentes et comparaison}
\begin{table}[H]
\centering
\caption{Comparaison avec des familles de solutions existantes}
\begin{tabular}{p{3.5cm}p{4.2cm}p{4.2cm}p{2.4cm}}
\toprule
\textbf{Type de solution} & \textbf{Fonctions principales} & \textbf{Limites} & \textbf{Position de Smart Firma} \\
\midrule
Applications météo & Prévisions, température, vent, pluie & Peu de contextualisation agronomique & Ajoute alertes agricoles et zones tunisiennes \\
Outils de diagnostic image & Détection de maladies végétales & Souvent isolés du reste du système agricole & Diagnostic intégré avec recommandations \\
Chatbots généralistes & Réponses conversationnelles & Contexte agricole limité & Prompt spécialisé agriculture tunisienne \\
Systèmes de recommandation & Conseils selon données tabulaires & Peu adaptés aux villes tunisiennes & Modèle localisé avec base sol et météo \\
\bottomrule
\end{tabular}
\end{table}

\subsection{Critique de l'existant}
Les solutions existantes sont souvent mono-fonctionnelles. Elles ne combinent pas toujours authentification, diagnostic, météo, recommandations, historiques et communication. Smart Firma répond à cette limite en intégrant plusieurs modules dans une architecture cohérente.

\section{Solution proposée}
La solution proposée est une application web client-serveur. Le client React consomme les APIs REST exposées par Django. Le backend orchestre les services métier, les APIs externes, les modèles Machine Learning et la base de données SQLite.

\begin{figure}[H]
\centering
\caption{Fonctionnement général de la solution}
\begin{lstlisting}[language=mermaid]
flowchart LR
    User[Agriculteur / Administrateur] --> UI[Interface React TypeScript]
    UI --> API[API Django REST]
    API --> DB[(SQLite)]
    API --> Weather[Open-Meteo]
    API --> News[NewsAPI]
    API --> OpenAI[OpenAI API]
    API --> RF[Random Forest crops]
    API --> CNN[CNN maladies]
\end{lstlisting}
\end{figure}

\section{Méthodologie de travail}
\subsection{Méthodes agiles}
Le projet est organisé selon une logique agile : les fonctionnalités sont découpées en modules livrables et testables. Chaque incrément apporte une valeur visible à l'utilisateur.

\subsection{SCRUM}
La démarche SCRUM est adaptée au projet avec les rôles suivants :
\begin{itemize}[noitemsep]
  \item Product Owner : priorisation des besoins et validation fonctionnelle ;
  \item Scrum Master : organisation du travail et suivi de l'avancement ;
  \item Equipe de développement : frontend, backend, intégration IA et tests ;
  \item Utilisateurs finaux : agriculteurs, administrateurs et techniciens.
\end{itemize}

\subsection{Cycle de développement}
Le cycle suit les étapes : analyse, conception, implémentation, test, correction et documentation. Les sprints sont détaillés dans le chapitre 3.

\chapter{Analyse et conception}

\section{Capture des besoins}
\subsection{Besoins fonctionnels}
\begin{longtable}{p{1.5cm}p{10cm}p{2cm}}
\caption{Besoins fonctionnels identifiés}\\
\toprule
\textbf{ID} & \textbf{Besoin} & \textbf{Module} \\
\midrule
BF1 & Créer un compte utilisateur avec nom, prénom, email, téléphone et mot de passe. & Auth \\
BF2 & Se connecter et conserver la session avec un token. & Auth \\
BF3 & Consulter et mettre à jour le profil. & Auth \\
BF4 & Consulter un tableau de bord utilisateur. & Frontend \\
BF5 & Consulter un tableau de bord administrateur avec nombre d'utilisateurs. & Admin \\
BF6 & Obtenir la météo actuelle et les prévisions. & Weather \\
BF7 & Recommander une culture selon localisation et données agro-climatiques. & Crops ML \\
BF8 & Personnaliser N, P, K, pH, température et humidité. & Crops ML \\
BF9 & Téléverser une image de plante pour diagnostic. & Analysis DL \\
BF10 & Obtenir top-3 prédictions, confiance et recommandations. & Analysis DL \\
BF11 & Echanger des messages internes. & Messaging \\
BF12 & Discuter avec un chatbot agricole. & Chatbot \\
BF13 & Consulter les actualités. & News \\
\bottomrule
\end{longtable}

\subsection{Besoins non fonctionnels}
\begin{itemize}[noitemsep]
  \item Sécurité : authentification par token DRF, permissions par endpoint, validation des entrées.
  \item Performance : chargement paresseux du CNN, cache LRU sur le géocodage des villes.
  \item Maintenabilité : applications Django séparées et services frontend isolés.
  \item Portabilité : rapport compatible Overleaf, application exécutable localement.
  \item Traçabilité : sauvegarde des analyses et conversations.
  \item Ergonomie : interfaces spécialisées par usage métier.
\end{itemize}

\subsection{Contraintes techniques}
Les contraintes détectées sont :
\begin{itemize}[noitemsep]
  \item backend Python/Django 4.2.7 ;
  \item frontend React 19 avec TypeScript ;
  \item base SQLite en environnement de développement ;
  \item dépendance à des APIs externes : Open-Meteo, NewsAPI, OpenAI ;
  \item modèle CNN volumineux \texttt{plant\_disease\_cnn.h5} ;
  \item modèles Random Forest sauvegardés en \texttt{.pkl}.
\end{itemize}

\section{Identification des acteurs}
\begin{table}[H]
\centering
\caption{Acteurs du système}
\begin{tabular}{p{3cm}p{10cm}}
\toprule
\textbf{Acteur} & \textbf{Description} \\
\midrule
Visiteur & Accède à l'accueil, à la connexion et à l'inscription. \\
Utilisateur & Accède aux services authentifiés : météo, recommandation, chatbot, analyse, messagerie. \\
Administrateur & Consulte le tableau de bord administrateur et la liste des utilisateurs via endpoint protégé. \\
API externe & Fournit météo, actualités ou réponses conversationnelles. \\
Modèle ML/DL & Produit des prédictions de culture ou de maladie. \\
\bottomrule
\end{tabular}
\end{table}

\section{Cas d'utilisation}
\subsection{Diagramme global}
\begin{figure}[H]
\centering
\caption{Diagramme global des cas d'utilisation}
\begin{lstlisting}[language=mermaid]
flowchart LR
    Visiteur((Visiteur))
    Utilisateur((Utilisateur))
    Admin((Administrateur))

    UC1[S'inscrire]
    UC2[Se connecter]
    UC3[Consulter meteo]
    UC4[Recommander culture]
    UC5[Analyser image plante]
    UC6[Utiliser chatbot]
    UC7[Gerer messages]
    UC8[Consulter dashboard admin]

    Visiteur --> UC1
    Visiteur --> UC2
    Utilisateur --> UC3
    Utilisateur --> UC4
    Utilisateur --> UC5
    Utilisateur --> UC6
    Utilisateur --> UC7
    Admin --> UC8
    Admin --> UC7
\end{lstlisting}
\end{figure}

\subsection{Descriptions textuelles}
\begin{longtable}{p{3cm}p{11cm}}
\caption{Descriptions des cas d'utilisation}\\
\toprule
\textbf{Cas} & \textbf{Description} \\
\midrule
Inscription & L'utilisateur renseigne ses informations. Le backend vérifie l'unicité de l'email et du nom d'utilisateur, compare les mots de passe et crée un compte. \\
Connexion & L'utilisateur fournit son nom d'utilisateur et son mot de passe. L'API authentifie et retourne un token stocké côté client. \\
Météo & L'utilisateur choisit une ville tunisienne. Le frontend appelle \texttt{/api/weather/current/} et \texttt{/api/weather/forecast/}. \\
Recommandation & L'utilisateur saisit une localisation et optionnellement des valeurs manuelles. Le backend récupère météo et sol, puis interroge Random Forest. \\
Analyse image & L'utilisateur envoie une image. L'API valide le fichier, prétraite l'image, appelle le CNN et sauvegarde le diagnostic. \\
Chatbot & L'utilisateur pose une question. Le backend envoie le contexte agricole et la question à l'API OpenAI puis sauvegarde la conversation. \\
Messagerie & L'utilisateur consulte ses messages reçus/envoyés, envoie un message, lit le détail ou supprime un message. \\
\bottomrule
\end{longtable}

\section{Modélisation UML}
\subsection{Diagramme de classes}
\begin{figure}[H]
\centering
\caption{Diagramme de classes principal}
\begin{lstlisting}[language=mermaid]
classDiagram
    class User {
      int id
      string username
      string email
      string role
      string phone
      image avatar
      datetime created_at
    }
    class Message {
      int id
      string subject
      text content
      boolean is_read
      datetime created_at
      datetime updated_at
    }
    class AnalysisLog {
      int id
      image image
      json result
      float confidence
      float processing_time
      datetime created_at
    }
    class ChatMessage {
      int id
      text message
      text response
      datetime created_at
    }
    User "1" --> "*" Message : sender
    User "1" --> "*" Message : recipient
    User "1" --> "*" AnalysisLog
    User "1" --> "*" ChatMessage
\end{lstlisting}
\end{figure}

\subsection{Diagramme de séquence : authentification}
\begin{figure}[H]
\centering
\caption{Séquence de connexion}
\begin{lstlisting}[language=mermaid]
sequenceDiagram
    actor U as Utilisateur
    participant F as Frontend React
    participant A as API Auth Django
    participant DB as SQLite
    U->>F: Saisir username/password
    F->>A: POST /api/auth/login/
    A->>DB: Verifier utilisateur
    DB-->>A: Utilisateur valide
    A-->>F: user + token
    F->>F: Stocker token localStorage
    F-->>U: Redirection dashboard
\end{lstlisting}
\end{figure}

\subsection{Diagramme de séquence : diagnostic image}
\begin{figure}[H]
\centering
\caption{Séquence d'analyse de maladie}
\begin{lstlisting}[language=mermaid]
sequenceDiagram
    actor U as Utilisateur
    participant F as Page Analysis
    participant API as API Analysis
    participant CNN as plant_disease_cnn.h5
    participant DB as SQLite
    U->>F: Selectionner image
    F->>API: POST /api/analysis/upload/
    API->>API: Valider type et taille
    API->>DB: Creer AnalysisLog
    API->>API: Resize 224x224, RGB, normalisation
    API->>CNN: Prediction
    CNN-->>API: Probabilites 38 classes
    API->>DB: Sauvegarder resultat JSON
    API-->>F: Top prediction + recommandations
    F-->>U: Afficher diagnostic
\end{lstlisting}
\end{figure}

\subsection{Diagramme d'activité : recommandation agricole}
\begin{figure}[H]
\centering
\caption{Activité du module de recommandation}
\begin{lstlisting}[language=mermaid]
flowchart TD
    A[Debut] --> B[Saisie localisation]
    B --> C[Geocodage Open-Meteo]
    C --> D{Meteo disponible ?}
    D -- Non --> E[Retour erreur 503]
    D -- Oui --> F[Charger donnees sol]
    F --> G{Mode manuel ?}
    G -- Oui --> H[Surcharger NPK, pH, temperature, humidite]
    G -- Non --> I[Utiliser donnees automatiques]
    H --> J[Construire DataFrame]
    I --> J
    J --> K[Predire avec RandomForest]
    K --> L[Calculer confiance et alternatives]
    L --> M[Generer conseils]
    M --> N[Retour JSON]
\end{lstlisting}
\end{figure}

\subsection{Diagramme de composants}
\begin{figure}[H]
\centering
\caption{Diagramme de composants}
\begin{lstlisting}[language=mermaid]
flowchart LR
    subgraph Frontend
        Pages[Pages React]
        Services[Services Axios]
        Context[AuthContext]
    end

    subgraph Backend
        Auth[authentication]
        Weather[weather]
        News[news]
        Analysis[analysis]
        Crops[crops]
        Chatbot[chatbot]
        Messaging[messaging]
    end

    subgraph Data
        SQLite[(db.sqlite3)]
        Media[(media files)]
        Models[(PKL/H5 models)]
    end

    Pages --> Services
    Context --> Services
    Services --> Auth
    Services --> Weather
    Services --> News
    Services --> Analysis
    Services --> Chatbot
    Services --> Messaging
    Pages --> Crops
    Auth --> SQLite
    Messaging --> SQLite
    Analysis --> SQLite
    Chatbot --> SQLite
    Analysis --> Media
    Analysis --> Models
    Crops --> Models
\end{lstlisting}
\end{figure}

\section{Architecture logicielle}
L'application utilise une architecture client-serveur de type trois tiers. Le frontend React constitue la couche présentation. Le backend Django REST constitue la couche métier et expose les APIs. La base SQLite, les fichiers médias et les fichiers modèles constituent la couche données.

\subsection{Architecture frontend}
Le frontend est organisé autour des dossiers :
\begin{itemize}[noitemsep]
  \item \texttt{pages} : pages routées par React Router ;
  \item \texttt{components} : composants réutilisables dont \texttt{Navbar}, \texttt{Footer} et \texttt{RecommendationForm} ;
  \item \texttt{services} : clients API Axios ;
  \item \texttt{contexts} : contexte d'authentification ;
  \item \texttt{styles} : feuilles CSS.
\end{itemize}

\subsection{Architecture backend}
Le backend est organisé en applications Django indépendantes :
\texttt{authentication}, \texttt{messaging}, \texttt{weather}, \texttt{news}, \texttt{analysis}, \texttt{chatbot}, \texttt{crops}. Cette séparation améliore la maintenabilité et limite le couplage.

\section{Base de données}
\subsection{MCD}
\begin{figure}[H]
\centering
\caption{MCD simplifié}
\begin{lstlisting}[language=mermaid]
erDiagram
    USER ||--o{ MESSAGE : sends
    USER ||--o{ MESSAGE : receives
    USER ||--o{ ANALYSIS_LOG : owns
    USER ||--o{ CHAT_MESSAGE : owns

    USER {
      int id
      string username
      string email
      string role
      string phone
      string avatar
    }
    MESSAGE {
      int id
      int sender_id
      int recipient_id
      string subject
      text content
      boolean is_read
    }
    ANALYSIS_LOG {
      int id
      int user_id
      string image
      json result
      float confidence
      float processing_time
    }
    CHAT_MESSAGE {
      int id
      int user_id
      text message
      text response
    }
\end{lstlisting}
\end{figure}

\subsection{Explication des entités}
\begin{itemize}[noitemsep]
  \item \texttt{User} : utilisateur personnalisé avec rôle, téléphone et avatar.
  \item \texttt{Message} : message interne entre deux utilisateurs.
  \item \texttt{AnalysisLog} : historique d'analyse image, résultat JSON et métriques d'inférence.
  \item \texttt{ChatMessage} : historique des échanges avec le chatbot.
\end{itemize}

\section{Gestion de projet SCRUM}
\subsection{Product backlog}
\begin{longtable}{p{1.2cm}p{8.5cm}p{2cm}p{2cm}}
\caption{Product backlog}\\
\toprule
\textbf{ID} & \textbf{User story} & \textbf{Priorité} & \textbf{Sprint} \\
\midrule
US1 & En tant que visiteur, je veux créer un compte. & Haute & S1 \\
US2 & En tant qu'utilisateur, je veux me connecter. & Haute & S1 \\
US3 & En tant qu'utilisateur, je veux consulter mon tableau de bord. & Haute & S1 \\
US4 & En tant qu'utilisateur, je veux consulter la météo agricole. & Haute & S2 \\
US5 & En tant qu'utilisateur, je veux consulter les actualités. & Moyenne & S2 \\
US6 & En tant qu'agriculteur, je veux obtenir une recommandation de culture. & Très haute & S2 \\
US7 & En tant qu'agriculteur, je veux personnaliser les valeurs du sol. & Haute & S3 \\
US8 & En tant qu'utilisateur, je veux diagnostiquer une maladie à partir d'une image. & Très haute & S3 \\
US9 & En tant qu'utilisateur, je veux recevoir des conseils de traitement. & Très haute & S3 \\
US10 & En tant qu'utilisateur, je veux discuter avec un chatbot. & Moyenne & S4 \\
US11 & En tant qu'utilisateur, je veux envoyer et recevoir des messages. & Moyenne & S4 \\
US12 & En tant qu'administrateur, je veux consulter un tableau de bord. & Moyenne & S4 \\
\bottomrule
\end{longtable}

\subsection{Planification des sprints}
\begin{table}[H]
\centering
\caption{Planification des sprints}
\begin{tabular}{p{2cm}p{5cm}p{7cm}}
\toprule
\textbf{Sprint} & \textbf{Objectif} & \textbf{Livrables} \\
\midrule
S1 & Socle applicatif & Structure React/Django, routage, authentification, dashboards. \\
S2 & Services contextuels & Météo, actualités, recommandation initiale de cultures. \\
S3 & Intelligence artificielle & Random Forest Tunisie, CNN maladies, upload image, recommandations. \\
S4 & Communication et consolidation & Messagerie, chatbot, dashboard admin, documentation et tests. \\
\bottomrule
\end{tabular}
\end{table}

\chapter{Réalisation}

\section{Environnement technique}
\begin{table}[H]
\centering
\caption{Technologies utilisées}
\begin{tabular}{p{4cm}p{9cm}}
\toprule
\textbf{Couche} & \textbf{Technologies} \\
\midrule
Frontend & React 19, TypeScript, React Router DOM, Axios, Lucide React, CSS. \\
Backend & Python, Django 4.2.7, Django REST Framework, django-cors-headers. \\
Base de données & SQLite en développement. \\
Machine Learning & Scikit-learn, RandomForestClassifier, Joblib, Pandas, NumPy. \\
Deep Learning & TensorFlow 2.17, Keras, Pillow. \\
Services externes & Open-Meteo, OpenWeatherMap, NewsAPI, OpenAI API. \\
\bottomrule
\end{tabular}
\end{table}

\section{Sprint 1 : socle applicatif et authentification}
\subsection{Objectifs}
Le premier sprint vise à mettre en place la structure globale du projet, l'authentification et les premières interfaces.

\subsection{Backlog du sprint}
\begin{itemize}[noitemsep]
  \item initialiser le backend Django ;
  \item créer le frontend React ;
  \item mettre en place les routes principales ;
  \item créer le modèle utilisateur personnalisé ;
  \item développer inscription, connexion, profil et déconnexion ;
  \item créer les dashboards utilisateur et administrateur.
\end{itemize}

\subsection{Réalisation}
Le backend contient l'application \texttt{authentication}. Le modèle \texttt{User} étend \texttt{AbstractUser} et ajoute les attributs \texttt{role}, \texttt{phone}, \texttt{avatar}, \texttt{created\_at} et \texttt{updated\_at}. Les serializers vérifient l'unicité de l'email et du nom d'utilisateur, ainsi que la concordance des mots de passe.

Côté frontend, le contexte \texttt{AuthContext} centralise l'état d'authentification. Le service \texttt{api.ts} ajoute automatiquement le header \texttt{Authorization: Token <token>} à chaque requête. En cas d'erreur 401, le token est supprimé et l'utilisateur est redirigé vers la page de connexion.

\subsection{Interfaces développées}
\begin{longtable}{p{3cm}p{11cm}}
\caption{Interfaces du sprint 1}\\
\toprule
\textbf{Page} & \textbf{Rôle et fonctionnement} \\
\midrule
Accueil & Présente Smart Firma, les services agricoles et les accès connexion/inscription ou dashboard selon l'état de session. \\
Login & Formulaire username/password. Après succès, redirection vers dashboard utilisateur ou admin. \\
Register & Formulaire d'inscription avec prénom, nom, username, email, téléphone optionnel et confirmation du mot de passe. \\
Dashboard & Affiche données utilisateur, zones agro-écologiques, statistiques agricoles et raccourcis vers les services. \\
DashboardAdmin & Affiche nombre d'utilisateurs, messages non lus et cartes d'administration. \\
\bottomrule
\end{longtable}

\subsection{Captures d'écran des interfaces du sprint}
Les captures d'écran suivantes documentent les principales interfaces réalisées pendant le sprint 1. Les fichiers doivent être placés dans le dossier \texttt{docs/rapport/figures/}. Si une image n'est pas encore disponible, le rapport affiche automatiquement un cadre réservé à la capture afin de conserver la cohérence de la mise en page.

\screenshotfigure{01_accueil.png}{Interface d'accueil de la plateforme Smart Firma}{fig:accueil}

\screenshotfigure{02_login.png}{Interface de connexion utilisateur}{fig:login}

\screenshotfigure{03_register.png}{Interface d'inscription utilisateur}{fig:register}

\screenshotfigure{04_dashboard_utilisateur.png}{Tableau de bord utilisateur}{fig:dashboard-user}

\screenshotfigure{05_dashboard_admin.png}{Tableau de bord administrateur}{fig:dashboard-admin}

\subsection{APIs implémentées}
\begin{longtable}{p{4.5cm}p{2cm}p{7cm}}
\caption{APIs du sprint 1}\\
\toprule
\textbf{Endpoint} & \textbf{Méthode} & \textbf{Description} \\
\midrule
\texttt{/api/auth/register/} & POST & Inscription et création de token. \\
\texttt{/api/auth/login/} & POST & Authentification et retour du token. \\
\texttt{/api/auth/logout/} & POST & Suppression du token utilisateur. \\
\texttt{/api/auth/profile/} & GET & Consultation du profil connecté. \\
\texttt{/api/auth/profile/update/} & PUT & Mise à jour partielle du profil. \\
\texttt{/api/auth/users/} & GET & Liste des utilisateurs, réservée aux administrateurs. \\
\bottomrule
\end{longtable}

\section{Sprint 2 : services agricoles contextuels}
\subsection{Objectifs}
Ce sprint ajoute les services météo, actualités et recommandation agricole initiale.

\subsection{Réalisation météo}
Le module \texttt{weather} interroge Open-Meteo en deux étapes : géocodage de la ville puis récupération des données météo. Il fournit la météo actuelle et une prévision jusqu'à cinq jours. Le backend prévoit un fallback vers OpenWeatherMap si une clé est configurée, puis des données simulées en dernier recours.

\subsection{Réalisation actualités}
Le module \texttt{news} utilise NewsAPI. Il retourne les articles avec titre, description, contenu, auteur, source, URL, image et date de publication.

\subsection{Réalisation recommandation}
Le module \texttt{crops} charge \texttt{crop\_model\_tunisia.pkl}. Il récupère la localisation, les données météo Open-Meteo et les données sol locales à partir d'un dictionnaire couvrant plusieurs villes tunisiennes. Il construit le vecteur \([N,P,K,temperature,humidity,ph,rainfall]\) puis appelle \texttt{predict} et \texttt{predict\_proba}.

\subsection{Interfaces développées}
\begin{longtable}{p{3cm}p{11cm}}
\caption{Interfaces du sprint 2}\\
\toprule
\textbf{Page} & \textbf{Description} \\
\midrule
Weather & Recherche par ville, affichage météo actuelle, prévisions 5 jours, zone agro-écologique et alertes agricoles. \\
RecommendationPage & Page dédiée au module Smart-Farm Tunisia et intégration du composant \texttt{RecommendationForm}. \\
RecommendationForm & Formulaire de localisation, mode manuel NPK/pH/température/humidité, affichage culture recommandée, alternatives, météo, sol et conseils. \\
\bottomrule
\end{longtable}

\subsection{Captures d'écran des interfaces du sprint}
\screenshotfigure{06_meteo.png}{Interface de consultation météo agricole}{fig:meteo}

\screenshotfigure{07_recommandation_formulaire.png}{Formulaire de recommandation agricole Smart-Farm Tunisia}{fig:recommandation-form}

\screenshotfigure{08_recommandation_resultat.png}{Résultat de recommandation de culture avec alternatives et conseils}{fig:recommandation-resultat}

\section{Sprint 3 : intelligence artificielle et diagnostic}
\subsection{Objectifs}
Ce sprint intègre les modèles Machine Learning et Deep Learning.

\subsection{Module Random Forest}
Le dataset tunisien contient 1710 observations et 11 cultures. Le modèle Random Forest tunisien est entraîné avec 200 arbres, profondeur maximale 15 et \texttt{min\_samples\_split=5}. Les résultats mesurés indiquent une validation croisée moyenne de 94,44 \% et une accuracy holdout de 95,61 \%.

\subsection{Module CNN}
Le modèle \texttt{plant\_disease\_cnn.h5} est un réseau Keras séquentiel pour 38 classes. Le backend charge le modèle paresseusement afin d'éviter un coût important à chaque requête. Le fichier image est validé, limité à 10 MB, converti en RGB, redimensionné à 224x224 et normalisé.

\subsection{Interface Analysis}
La page \texttt{Analysis.tsx} permet de sélectionner une image, visualiser l'aperçu, lancer l'analyse, afficher la confiance, le top diagnostic, les alternatives, les recommandations de traitement/prévention, le temps d'inférence et l'historique local.

\subsection{Captures d'écran des interfaces du sprint}
\screenshotfigure{09_analyse_upload.png}{Interface de téléversement d'une image de plante}{fig:analyse-upload}

\screenshotfigure{10_analyse_resultat.png}{Résultat du diagnostic de maladie avec confiance et recommandations}{fig:analyse-resultat}

\section{Partie Machine Learning et IA}
\subsection{Introduction}
La solution utilise le Machine Learning supervisé pour la recommandation de cultures et le Deep Learning pour la classification d'images de maladies végétales. Ces deux modules répondent à des besoins différents : décision agro-climatique et diagnostic visuel.

\subsection{Méthodologie CRISP-DM}
\begin{table}[H]
\centering
\caption{Application de CRISP-DM au projet}
\begin{tabular}{p{3.2cm}p{10cm}}
\toprule
\textbf{Phase} & \textbf{Application} \\
\midrule
Business understanding & Aider les agriculteurs tunisiens à choisir une culture et diagnostiquer des maladies. \\
Data understanding & Analyse des datasets CSV, classes agricoles, variables NPK, météo, pH et images de feuilles. \\
Data preparation & Séparation X/y pour tabulaire, génération dataset Tunisie, redimensionnement et normalisation des images. \\
Modeling & Random Forest pour cultures, CNN Keras pour maladies. \\
Evaluation & Cross-validation, holdout accuracy, confiance de prédiction, top-3 CNN. \\
Deployment & Endpoints REST \texttt{/crops/predict\_crop/} et \texttt{/api/analysis/upload/}. \\
\bottomrule
\end{tabular}
\end{table}

\subsection{Datasets}
Le dataset général \texttt{Crop\_recommendation.csv} contient 2200 lignes, 22 classes équilibrées et 7 variables numériques. Le dataset \texttt{crop\_recommendation\_tunisia.csv} contient 1710 lignes et 11 cultures tunisiennes. Le dataset image n'est pas fourni sous forme de dossier d'entraînement dans le workspace ; seules des images médias d'analyses et le modèle CNN entraîné sont présents. Le code indique que les classes correspondent au dataset public \texttt{vipoooool/new-plant-diseases-dataset}.

\subsection{Nettoyage, valeurs manquantes et outliers}
Le code d'entraînement ne contient pas de traitement explicite des valeurs manquantes ou des outliers. Les données tabulaires sont utilisées directement après lecture CSV. Pour un pipeline de production, il est recommandé d'ajouter :
\begin{itemize}[noitemsep]
  \item contrôle des valeurs nulles ;
  \item validation des bornes physiques N, P, K, pH, température, humidité et pluie ;
  \item détection des valeurs extrêmes ;
  \item journalisation des données rejetées.
\end{itemize}

\subsection{Prétraitement et normalisation}
Le Random Forest ne nécessite pas de normalisation stricte. Le CNN applique une normalisation pixel par division par 255.0. Les images sont converties en tenseurs \texttt{float32}.

\subsection{Analyse exploratoire}
Les statistiques extraites du dataset tunisien montrent des pH principalement alcalins, une pluviométrie moyenne de 335,48 mm et une température moyenne de 22,25 °C. Les cultures sont volontairement distribuées selon les régions Nord, Centre et Sud.

\subsection{Modélisation et hyperparamètres}
\begin{longtable}{p{4cm}p{5cm}p{5cm}}
\caption{Modèles IA utilisés}\\
\toprule
\textbf{Modèle} & \textbf{Hyperparamètres} & \textbf{Usage} \\
\midrule
Random Forest général & 100 arbres, \texttt{random\_state=42} & Recommandation générale de cultures. \\
Random Forest Tunisie & 200 arbres, \texttt{max\_depth=15}, \texttt{min\_samples\_split=5} & Recommandation adaptée à la Tunisie. \\
CNN Keras & Conv2D 32/64/128, Dense 128, Dropout 0.5, Softmax 38 & Classification d'images de maladies. \\
\bottomrule
\end{longtable}

\subsection{Evaluation}
\begin{table}[H]
\centering
\caption{Résultats d'évaluation disponibles}
\begin{tabular}{lrrr}
\toprule
\textbf{Modèle} & \textbf{CV moyenne} & \textbf{Ecart-type} & \textbf{Accuracy holdout} \\
\midrule
Random Forest général & 0.9945 & 0.0034 & 0.9955 \\
Random Forest Tunisie & 0.9444 & 0.0256 & 0.9561 \\
\bottomrule
\end{tabular}
\end{table}

Les métriques precision, recall et F1-score ne sont pas persistées dans le projet, mais elles sont recommandées pour compléter l'évaluation. Pour le CNN, le code retourne une confiance et un top-3 mais ne contient pas de rapport de test indépendant.

\subsection{Déploiement IA}
Le déploiement est effectué directement dans le backend Django. Les modèles sont stockés dans le dossier \texttt{backend}. Le modèle Random Forest est chargé au chargement du module \texttt{crops.views}. Le CNN est chargé uniquement lors de la première demande d'analyse.

\section{Sprint 4 : chatbot, messagerie et consolidation}
\subsection{Objectifs}
Le quatrième sprint consolide l'expérience utilisateur avec communication et assistance intelligente.

\subsection{Messagerie}
Le module \texttt{messaging} permet de consulter la boîte de réception, les messages envoyés, le détail d'un message, l'envoi d'un message, la suppression et le comptage des messages non lus. L'accès au détail vérifie que l'utilisateur est expéditeur ou destinataire.

\subsection{Chatbot}
Le module \texttt{chatbot} utilise un prompt système orienté agriculture tunisienne. Les messages sont enregistrés dans \texttt{ChatMessage}. En cas d'erreur côté frontend, des réponses de secours sont prévues dans la page \texttt{Chatbot.tsx}.

\subsection{Captures d'écran des interfaces du sprint}
\screenshotfigure{11_chatbot.png}{Interface du chatbot agricole spécialisé}{fig:chatbot}

\screenshotfigure{12_messagerie.png}{Interface ou logique de messagerie interne}{fig:messagerie}

\section{Sécurité}
\subsection{Authentification et autorisation}
Le projet utilise \texttt{TokenAuthentication} et \texttt{SessionAuthentication}. Il n'utilise pas JWT ni OAuth dans le code source observé. Les permissions DRF globales imposent \texttt{IsAuthenticated}, avec exceptions explicites pour l'inscription et la connexion.

\subsection{Validation}
Les validations principales sont :
\begin{itemize}[noitemsep]
  \item unicité email et username ;
  \item longueur minimale du mot de passe ;
  \item confirmation du mot de passe ;
  \item type MIME image ;
  \item taille image maximale de 10 MB ;
  \item contrôle d'accès aux messages et analyses.
\end{itemize}

\subsection{Middleware et CORS}
Le middleware \texttt{corsheaders.middleware.CorsMiddleware} est activé. Les origines CORS sont configurables via variable d'environnement.

\subsection{Gestion des erreurs}
Les endpoints retournent des codes HTTP adaptés : 400 pour données invalides, 403 pour accès interdit, 503 pour indisponibilité du modèle ou de la météo, 500 pour erreurs internes.

\section{DevOps}
Aucun fichier Docker, pipeline CI/CD, configuration Kubernetes ou monitoring dédié n'a été détecté dans le workspace. Le déploiement actuel est donc local, avec \texttt{python manage.py runserver} pour Django et \texttt{npm start} pour React. Pour une mise en production, il est recommandé d'ajouter Docker, PostgreSQL, Nginx, HTTPS, logs applicatifs et une chaîne CI/CD.

\section{Tests et résultats}
Des fichiers \texttt{tests.py} existent dans les applications Django, mais ils ne contiennent pas de stratégie de tests détaillée dans le code inspecté. Les tests recommandés sont :
\begin{itemize}[noitemsep]
  \item tests unitaires serializers ;
  \item tests API authentification ;
  \item tests de permissions messages et analyses ;
  \item tests de validation image ;
  \item tests d'intégration frontend/backend ;
  \item tests ML avec matrice de confusion et rapports de classification.
\end{itemize}

\chapter{Conclusion générale}

Le projet Smart Firma constitue une plateforme web intelligente orientée agriculture tunisienne. L'analyse du code source montre une architecture claire, structurée autour d'un frontend React/TypeScript et d'un backend Django REST Framework. Les fonctionnalités principales couvrent l'authentification, les tableaux de bord, la météo agricole, les actualités, la messagerie, le chatbot, la recommandation de cultures et le diagnostic de maladies végétales.

Sur le plan technique, la solution exploite des services externes et des modèles d'intelligence artificielle. Le module de recommandation de cultures repose sur un Random Forest entraîné avec des données agro-climatiques. Le modèle tunisien obtient une accuracy holdout de 95,61 \%, ce qui constitue un résultat encourageant pour un prototype. Le module de diagnostic utilise un CNN Keras capable de classifier 38 classes de plantes et maladies à partir d'images redimensionnées en 224x224.

\section*{Résultats obtenus}
\addcontentsline{toc}{section}{Résultats obtenus}
\begin{itemize}[noitemsep]
  \item architecture client-serveur fonctionnelle ;
  \item APIs REST modulaires ;
  \item authentification par token ;
  \item intégration météo et actualités ;
  \item recommandation agricole adaptée aux zones tunisiennes ;
  \item diagnostic image par Deep Learning ;
  \item chatbot spécialisé agriculture tunisienne ;
  \item documentation académique et technique générée.
\end{itemize}

\section*{Limites}
\addcontentsline{toc}{section}{Limites}
Les principales limites sont l'absence de Docker/CI-CD, l'utilisation de SQLite pour le développement, le manque de tests automatisés complets, l'absence d'un rapport d'entraînement détaillé pour le CNN et le caractère synthétique du dataset tunisien. De plus, le projet utilise TokenAuthentication mais ne contient ni JWT ni OAuth.

\section*{Perspectives}
\addcontentsline{toc}{section}{Perspectives}
Les perspectives d'amélioration sont :
\begin{itemize}[noitemsep]
  \item migration vers PostgreSQL ;
  \item ajout de Docker et d'une pipeline CI/CD ;
  \item enrichissement du dataset tunisien avec des données terrain ;
  \item ajout de rapports ML complets : confusion matrix, precision, recall, F1-score ;
  \item amélioration du CNN par transfert learning ;
  \item ajout de capteurs IoT et données parcellaire ;
  \item gestion avancée des rôles et permissions ;
  \item déploiement cloud sécurisé avec HTTPS et monitoring.
\end{itemize}

\appendix
\chapter{Annexes}

\section{Endpoints API}
\begin{longtable}{p{5cm}p{2cm}p{7cm}}
\caption{Documentation synthétique des APIs}\\
\toprule
\textbf{Endpoint} & \textbf{Méthode} & \textbf{Description} \\
\midrule
\texttt{/api/auth/register/} & POST & Inscription utilisateur. \\
\texttt{/api/auth/login/} & POST & Connexion et retour token. \\
\texttt{/api/auth/logout/} & POST & Déconnexion et suppression token. \\
\texttt{/api/auth/profile/} & GET & Profil connecté. \\
\texttt{/api/auth/profile/update/} & PUT & Mise à jour profil. \\
\texttt{/api/auth/users/} & GET & Liste utilisateurs admin. \\
\texttt{/api/weather/current/} & GET & Météo actuelle. \\
\texttt{/api/weather/forecast/} & GET & Prévisions météo. \\
\texttt{/api/news/} & GET & Actualités filtrées. \\
\texttt{/api/news/categories/} & GET & Catégories d'actualités. \\
\texttt{/api/analysis/upload/} & POST & Upload image et diagnostic. \\
\texttt{/api/analysis/history/} & GET & Historique analyses. \\
\texttt{/api/analysis/<id>/} & GET & Détail analyse. \\
\texttt{/api/chatbot/chat/} & POST & Message chatbot. \\
\texttt{/api/chatbot/history/} & GET & Historique chatbot. \\
\texttt{/api/chatbot/clear/} & DELETE & Effacer historique chatbot. \\
\texttt{/api/messaging/inbox/} & GET & Messages reçus. \\
\texttt{/api/messaging/sent/} & GET & Messages envoyés. \\
\texttt{/api/messaging/send/} & POST & Envoyer message. \\
\texttt{/api/messaging/<id>/} & GET & Détail message. \\
\texttt{/api/messaging/<id>/delete/} & DELETE & Supprimer message. \\
\texttt{/api/messaging/unread/count/} & GET & Compteur non lus. \\
\texttt{/crops/predict\_crop/} & GET & Recommandation culture. \\
\texttt{/crops/locations/} & GET & Liste villes disponibles. \\
\bottomrule
\end{longtable}

\section{Commandes utiles}
\begin{lstlisting}[language=bash]
# Backend
cd backend
pip install -r requirements.txt
python manage.py migrate
python manage.py runserver

# Frontend
cd frontend
npm install
npm start

# Compilation Overleaf
# Ouvrir docs/rapport/main.tex puis compiler avec biber si necessaire.
\end{lstlisting}

\section{Extrait de configuration API frontend}
\begin{lstlisting}[language=JavaScript]
api.interceptors.request.use((config) => {
  const token = localStorage.getItem('token');
  if (token) {
    config.headers.Authorization = `Token ${token}`;
  }
  return config;
});
\end{lstlisting}

\section{Classes CNN}
Le CNN couvre 38 classes : pommier, myrtillier, cerisier, maïs, vigne, orange, pêcher, poivron, pomme de terre, framboisier, soja, courge, fraisier et tomate, avec états sains et maladies telles que tavelure, black rot, rouille, mildiou, alternariose, septoriose, oïdium, acariens, virus TYLCV et mosaïque.


\printbibliography[title={Bibliographie et webographie}]

\end{document}
